\subsection{The CM field and explicit nonarithmetical finite shadows}

\begin{statement}{D.25}{Proved}{A finite CM entanglement field and nonarithmetical GT0 shadows}
Put $r_0=(3/4)^{1/4}$ and $M'=\mathbb Q(i,r_0)$. There are exact equalities
\[
 \mathbb Q(E_3[4])=M'=\mathbb Q\bigl(x(E_4[3])\bigr),\qquad
 M'\mathbb Q^{\rm ab}=M\mathbb Q^{\rm ab},                  \tag{D.25.1}
\]
where $E_3:y^2=x^3+1$, $E_4:y^2=x^3-x$, and $M$ is D.16's field.
The projective $E_3[12]$ representation recovers the full $E_3[4]$
representation: choose its unique global sign fixing the rational
point $(0,1)$ of order three. Thus the extra relation of D.16 is
already detected by $T\bmod12$ and $W\bmod3$, with this sign rule.

Here is a finite refinement of the matrix problem, including
generation and the two hexagons. Write
\[
 B=\langle s,r:s^2=r^3=1\rangle,\quad
 Q=\mathrm{SL}_2(\mathbb Z/12)/\{\pm I\},\quad
 N=\ker(B\to Q),\quad G_j=B/[N,N]N^j.
\]
Use $s\mapsto\left(\begin{smallmatrix}0&1\\-1&0\end{smallmatrix}\right)$
and $r\mapsto\left(\begin{smallmatrix}0&-1\\1&-1\end{smallmatrix}\right)$
projectively. Then $|Q|=576$, $N$ is free of rank 97, and
$|G_j|=576j^{97}$ for $j=2,3,6$.
Among the 48 pairs of matrix conjugators with multiplier one and
$V_{12}=I$,24 satisfy the CM relation and 24 violate it.
Exact computation in $G_2$ and $G_3$ gives
\[
\begin{array}{c|rr}
\text{condition}&\text{CM-compatible}&\text{extra CM coset}\\\hline
\text{matrix conjugators}&24&24\\
\text{cusp relation in }G_2&12&12\\
\text{cusp relation in }G_3&24&24\\
\text{automorphisms of }G_6&12&12.
\end{array}                                               \tag{D.25.2}
\]
Each surviving row has an explicit certificate whose induced
matrix on $N^{\rm ab}/p$ has rank 97 for $p=2,3$.
The twelve extra pairs yield twelve distinct nonarithmetical
finite $\widehat{\rm GT}_0$ shadows for the pullback of
$[N,N]N^6$ to $B_3$. One representative has
\[
 A_{12}=I,\qquad
 \beta_{12}=\begin{pmatrix}9&8\\4&9\end{pmatrix},\qquad
 V_{12}=I.                                                \tag{D.25.3}
\]
Here ``nonarithmetical'' means that no element of $G_{\mathbb Q}$
has that finite shadow. D.27 proves that these twelve
\emph{selected} shadows have no genuine full profinite lift.
It also constructs a different shadow with the same CM
pair (D.25.3) that survives both class-two refinements;
genuine realization of that CM pair remains undecided.
\end{statement}
\begin{proof}
For a nonzero two-torsion point of $E_3$ with $x$-coordinate
$e\in\{-1,-\zeta_3,-\zeta_3^2\}$, the coordinates of its halves
satisfy $x=e(1\pm\sqrt3)$. Their $y$-coordinates have squares
$-9\pm6\sqrt3$. For the plus sign choose $y$ with
\[
 y=\sqrt3(\sqrt3-1)r_0.
\]
The other square root is obtained using the product
$3\sqrt{-3}$, up to sign. These formulas prove
$\mathbb Q(E_3[4])=\mathbb Q(i,r_0)$ in both directions.
For $E_4$ the three-division polynomial is
$3x^4-6x^2-1$. Choose two roots $a,b$ with
$a^2=1+2/\sqrt3$, $b^2=1-2/\sqrt3$ and $ab=i/\sqrt3$.
Then $w=a+ib$ satisfies $w^4=4/3$; conversely
\[
 a=(w+2/w)/2,\qquad b=(w-2/w)/(2i).
\]
The splitting field is therefore the same $M'$. Since
$r_0=3^{1/4}/\sqrt2$, adjoining the abelian field
$\mathbb Q(\sqrt2)$ proves the second equality in (D.25.1).

The single global sign in this comparison is essential.
For $P\in E_3[4]$ and $Q_0=(0,1)\in E_3[3]$, the addition law gives
\[
 y(P)=\frac{x(P)^2}{4}
       \bigl(x(P-Q_0)-x(P+Q_0)\bigr).
\]
The points on the right belong to $E_3[12]$, proving directly
that its projective coordinate field contains the full
four-torsion field. Equivalently every Galois matrix on $E_3[3]$
fixes $Q_0$, so exactly one of the two simultaneous signs on
$E_3[12]$ has that property. The separate projective
$E_3[4]$ field is only $\mathbb Q(\zeta_{12})$; the displayed
addition formula is the precise map supplying the missing sign.

We describe the finite calculation so that its certificates
can be verified without enumerating $576\cdot6^{97}$ elements.
Let $F=\ker(B\to S_3)$ and $H=q(F)\subset Q$; then $|H|=96$.
The graph $N\backslash\mathcal T$ of the Bass--Serre tree of
$C_2*C_3$ has edges $Q$ and vertices
$Q/\langle s\rangle\amalg Q/\langle r\rangle$.
There are 576 edges and 480 vertices. A spanning tree gives97
integral fundamental cycles, a basis of $N^{\rm ab}$.
Write $e_q$ for the oriented edge from its $s$-vertex to its
$r$-vertex. The path cocycle into $\mathbb Z[Q]$ is
\[
 P(s)=0,\quad P(r)=e_1-e_r,\quad
 P(gh)=P(g)+q(g)P(h).                                    \tag{D.25.4}
\]
Its restriction identifies $N^{\rm ab}$ with the cycle group.
Reduction modulo $j$ realizes $G_j$ faithfully as the image
of $B$ in $(\mathbb Z/j)[Q]\rtimes Q$.

For $A\in C_H(r)$ and $\beta\in C_H(s)$ choose the recorded
paths $P_A,P_\beta$. There are respectively six and eight such
choices. A lift of the cusp-preserving relation
\[
 \beta s\beta^{-1}A rA^{-1}=sr
\]
exists modulo $p$ exactly when
\[
 (1-s)P_\beta+
 s\bigl((1-r)P_A+(A-1)P(r)\bigr)
 \in (1-s)C_p+s(1-r)C_p,\quad
 C_p=H_1(N\backslash\mathcal T,\mathbb F_p).               \tag{D.25.5}
\]
This is an explicit linear system. The right side has rank 96
at $p=2$ and 97 at $p=3$. Its solution counts are those in
(D.25.2). For every surviving pair the code gives corrections
$a,b\in C_p$ and checks generation as follows. The induced
equivariant map on edges sends
\[
 e_q\longmapsto q\bigl(P_A+a-P_\beta-b+e_A\bigr).
\]
Its restriction to the 97 fundamental cycles has rank 97 in
each characteristic. The homomorphism already induces the
identity on $Q$, so it is an automorphism of $G_p$.
Chinese remainders combine the two recorded corrections into
single conjugators in $G_6$, giving an automorphism there.
The code also multiplies the resulting group elements in
the semidirect model and checks both hexagon identities
directly. All spanning-tree choices, matrix pairs, cycle
coefficients and ranks are reproducible from
\texttt{d\_cm\_shadow.py}; its JSON retains the certificates.

The CM test in these coordinates is particularly small.
Use the unique representative of an element of $H$ whose
modulo-four image lies in the four-group generated by
$U^2,W^2$. For every one of the six $A$ choices, $A_4=I$.
The rational order-three line is $\ker(R-I)$ modulo 3, spanned
by $(1,-1)$. If $A_3^{-1}$ acts on this line by $-1$, the
global-sign normalization changes $T_4=I$ to $-I$, and the
restriction to $M'$ sends $r_0$ to $-r_0$; otherwise it fixes
$r_0$. On the other factor $\beta_3^{-1}$ is in
$\{\pm I,\pm S\}$; its class is nontrivial exactly for
$\pm S$, again sending $r_0$ to $-r_0$. Comparing these two
bits gives the 24/24 split. In (D.25.3) the bits are zero and
one, respectively.

The conjugator bits are invariants of the finite action,
not arbitrary choices in a matrix centralizer. The same
linear computation of (D.25.4) gives
\[
 q\bigl(C_{G_2}(s)\cap\operatorname{im}F\bigr)=1,
 \qquad
 q\bigl(C_{G_3}(r)\cap\operatorname{im}F\bigr)=1.          \tag{D.25.6}
\]
For example the first test asks whether
$(1-s)P_\beta$ lies in $(1-s)C_2$; only $q(\beta)=1$
survives. The second uses
$(1-r)P_A+(A-1)P(r)$ and $(1-r)C_3$, again with only the
identity surviving. Thus $G_6$ uniquely determines the
reductions of both normalized conjugators in $Q$.

Finally we give the exact passage to finite GT shadows.
Put $u=(sr)^{-1}$, $v=sus$, $x=u^2$, $y=v^2$; these are
the images of the standard braid generators and their
squares under $B_3\to B_3/Z(B_3)=B$.
Let $\alpha$ be one of the finite automorphisms just constructed.
It fixes $u$ and satisfies
$\alpha(s)=\beta s\beta^{-1}=sf$, where
$f=s\beta s\beta^{-1}$. Choose a lift of $\beta$ to the
discrete free group $F=\langle x,y\rangle$ and write its
abelianization as $a_0x+b_0y$. Since $s$ swaps $x,y$,
$f^{\rm ab}=(b_0-a_0)(x-y)$.
Set $k=b_0-a_0$, $\beta'=x^k\beta$, and
$f'=s\beta's(\beta')^{-1}$. Then $f'\in[F,F]$.
The gauge $\alpha'=\operatorname{Int}(x^k)\alpha$ still
fixes $u$ and satisfies
\[
 \alpha'(s)=sf',\qquad\alpha'(r)=(f')^{-1}r,
 \qquad\alpha'(v)=(f')^{-1}vf'.                           \tag{D.25.7}
\]
The last two equalities, with $r=uv$ and $s=uvu$, are
exactly the two $m=0$ hexagons (2.1)--(2.2) of
\cite[Definitions2.1 and 2.3]{dBDHP2026}.
The pullback $\widetilde N$ of $[N,N]N^6$ to $B_3$ is
normal, finite index, and contained in $PB_3$, since
$G_6\to S_3$ is the standard permutation quotient.
Here $\widetilde N_{\rm ord}=36$: the cusp $sr$ has orders
24,36,72 in $G_2,G_3,G_6$, respectively, its square has
order 36 in $G_6$, and the braid center is killed.
Thus $m=0$ gives the required unit multiplier one,
$f'$ has the required commutator image, and $\alpha'$ being
an automorphism gives surjectivity. This proves that
$(0,f'\bmod\widetilde N)$ is a finite GT shadow.

The gauge multiplies $A,\beta,V$ on the left by the same
$x^k$, so it fixes $T=A^{-1}V$ and $W=\beta^{-1}V$.
Equation(D.25.6) makes their reductions recoverable from
the finite shadow. The twelve different extra pairs have
different such coordinates, hence give distinct shadows.
No Galois element can give one of them: its two restrictions
to the single field $M'$ in (D.25.1) must agree. This proves
nonarithmeticality without assuming that an arbitrary
finite shadow lifts to $\widehat{\rm GT}_0$.
The exact all-finer-quotient test is D.22. The later
calculation D.27 rejects these specific selected shadows
as genuine full profinite actions while retaining an
alternate homology correction for the same extra CM pair.
\end{proof}
